\documentclass[11pt, a4paper]{article}

\usepackage[T1]{fontenc}
\usepackage[utf8]{inputenc}
\usepackage{tgpagella}
\usepackage[spanish, es-noshorthands]{babel}
\usepackage{amsmath, amssymb, bm}
\usepackage{tabularx}
\usepackage{booktabs}
\usepackage[most]{tcolorbox}
\usepackage[margin=2.5cm]{geometry}
\usepackage{ragged2e}
\usepackage{fancyhdr}
\usepackage[american]{circuitikz}

\pagestyle{fancy}
\fancyhf{}
\setlength{\headheight}{16pt}
\lhead{\textbf{UCB Sede Tarija} | Ing. Mecatrónica}
\chead{}
\rhead{IMT-121 Circuitos Electrónicos I | Guía Simulación RC}
\lfoot{Prof: M.Sc. Ing. Bernardo Quiroga Turdera}
\rfoot{Página \thepage}
\renewcommand{\headrulewidth}{0.4pt}

\definecolor{ucbblue}{RGB}{0, 51, 102}

\begin{document}

\begin{tcolorbox}[colback=ucbblue!5!white, colframe=ucbblue, title=\textbf{Guía de Simulación: Análisis Transitorio RC en LTspice}]
    \centering \textbf{Validación de Modelos Dinámicos en el Dominio del Tiempo}
\end{tcolorbox}

\section*{1. Objetivo}
Validar las predicciones analíticas (condiciones de frontera, inercia temporal $\tau$ y evolución) de un circuito RC de primer orden mediante el comando de simulación temporal \texttt{.tran} de LTspice.

\section*{2. Topología y Representación de la Conmutación}
Para simular el cierre de un interruptor en el tiempo sin introducir componentes complejos, modelaremos la excitación usando una fuente de tipo \texttt{PULSE}. El pulso arrancará en $0\,\text{V}$, manteniendo el circuito relajado, y subirá abruptamente a $9\,\text{V}$ en el instante $t_s = 10\,\text{ms}$.

\begin{center}
\begin{circuitikz}[scale=1, transform shape]
    \draw (0,0) to[V, l={\textbf{V1 (PULSE)}}, invert] (0,3) 
          to[R, l={$R_1=3\,\text{k}\Omega$}] (4,3);
    \draw (4,3) to[short, *-o] (5.5,3) node[above]{$a$};
    \draw (4,3) to[R, l_={$R_2=6\,\text{k}\Omega$}, invert] (4,0);
    \draw (5.5,3) to[C, l={$C=100\,\mu\text{F}$}, v={$v_C$}] (5.5,0);
    \draw (0,0) -- (5.5,0) node[ground]{};
    \node[anchor=east] at (-1, 2) {\small $0 \to 9\,\text{V}$ en $t=10\,\text{ms}$};
\end{circuitikz}
\end{center}

\section*{3. Paso a Paso en LTspice (Ruta Verificada)}

\begin{enumerate}
    \item \textbf{Dibujo del Esquema:} Arme el circuito conectando la fuente \texttt{V1}, las resistencias y el capacitor tal como indica el diagrama. Defina la referencia de tierra.
    \item \textbf{Etiquetas (Labels):} Utilice la herramienta \textit{Label Net} para nombrar el nodo de entrada como \texttt{in} (positivo de V1) y el nodo superior del capacitor como \texttt{a}.
    \item \textbf{Configuración de Fuente:} Haga clic derecho sobre la fuente \texttt{V1} y seleccione \textbf{Advanced}. Elija el formato \texttt{PULSE} y configure los siguientes parámetros:
    \begin{itemize}
        \item \texttt{Vinitial}: 0
        \item \texttt{Von}: 9
        \item \texttt{Tdelay}: 10m \textit{(El pulso arranca a los 10 ms)}
        \item \texttt{Trise / Tfall}: 100u \textit{(Flanco muy rápido, simula un switch ideal)}
        \item \texttt{Ton}: 2 \textit{(El pulso se mantiene en ON más tiempo del necesario)}
        \item \texttt{Tperiod}: 4
    \end{itemize}
    \item \textbf{Directiva de Simulación:} Inserte una directiva \texttt{.tran}. Para observar $\sim 5\tau$ (donde $\tau=0.2\,\text{s}$), ajustamos el \textit{Stop Time} en \texttt{1.1s} y un \textit{Maximum Timestep} de \texttt{1m}. El comando final que debe aparecer en pantalla es: \textbf{\texttt{.tran 0 1.1 0 1m}}.
    \item \textbf{Ejecución:} Corra la simulación (Run). \textbf{NO} seleccione opciones de inicialización forzada (\texttt{uic} o \texttt{.ic}). Queremos que el simulador halle el punto de operación natural (que será $0\,\text{V}$) antes de los 10 ms.
    \item \textbf{Visualización:} Utilice las puntas de prueba para observar \texttt{V(in)} y \texttt{V(a)}.
\end{enumerate}

\section*{4. Criterios de Troubleshooting (¿Qué hacer si no funciona?)}
\begin{itemize}
    \item \textit{La gráfica de V(a) es una línea plana en 6V:} Su fuente \texttt{V1} sigue configurada como DC $9\,\text{V}$ permanente. LTspice inicializó el capacitor ya cargado. Cambie la fuente a \texttt{PULSE}.
    \item \textit{La gráfica se corta antes de estabilizarse:} Su \textit{Stop Time} es muy corto. Debe ser al menos $t_s + 5\tau$.
    \item \textit{La curva baja en lugar de subir:} Reemplazó los nodos o midió el voltaje al revés. Verifique que la tierra esté en la parte inferior del circuito.
\end{itemize}

\end{document}
